% =====================================================================
% Banco diversificado --- Teorema da Superposicao
% 6 N1 + 9 N2 + 10 N3 + 8 N4 = 33 questoes
% Cada questao pratica possui circuito proprio; nao ha reutilizacao de uma
% mesma figura como substituto generico para enunciados diferentes.
% =====================================================================
\subsubsection*{Banco complementar --- Teorema da Superposição}
\begin{atencao}[Banco visual diversificado]
Nos exercícios práticos, o circuito faz parte do enunciado. Topologias, valores,
fontes e grandezas de interesse variam de uma questão para outra. A repetição é
usada apenas quando pedagogicamente necessária para comparar métodos.
\end{atencao}
% N1 --- Superposicao: 6 questoes visuais e distintas
\blocofixacao

\begin{questao}[1]{}
No circuito, deseja-se calcular a contribuição de $E_1$ para a corrente $I$ em $R_L$. Indique como as outras fontes devem ser desativadas.
\circuitoq{\begin{circuitikz}[american,scale=.78,transform shape]
\draw (0,0) to[\fonteV,l^={$E_1$}] (0,3.6) to[R,l={$R_1$}] (3.2,3.6) coordinate(a);
\draw (a) to[R,l_={$R_L$},i>^={$I$}] (3.2,0)--(0,0);
\draw (6.4,0) to[isource,l_={$I_2$}] (6.4,3.6)--(a); \draw (6.4,0)--(3.2,0);
\draw (9.0,0) to[\fonteV,l_={$E_3$}] (9.0,3.6) to[R,l_={$R_3$}] (a); \draw (9.0,0)--(6.4,0);
\end{circuitikz}}
\resposta{$I_2$ aberto; $E_3$ em curto; $E_1$ permanece ativa}
\resolucao{Na superposição, apenas uma fonte independente fica ativa por vez. Fonte de corrente ideal zerada vira circuito aberto; fonte de tensão ideal zerada vira curto-circuito.}
\end{questao}

\begin{questao}[1]{}
Considere a referência de corrente indicada. Determine somente a contribuição de $E_1=\SI{16}{\volt}$.
\circuitoq{\begin{circuitikz}[american,scale=.9,transform shape]
\draw (0,0) to[\fonteV,l^={$\SI{16}{\volt}$}] (0,3.2) to[R,l={$\SI{8}{\ohm}$},i>^={$I$}] (4.6,3.2)
 to[\fonteV,l_={$\SI{4}{\volt}$},invert] (4.6,0)--(0,0);
\end{circuitikz}}
\resposta{$I_{E_1}=\SI{2}{\ampere}$}
\resolucao{Zerando a fonte de $\SI{4}{\volt}$, ela vira um curto. Assim $I_{E_1}=16/8=\SI{2}{\ampere}$ no sentido da referência.}
\end{questao}

\begin{questao}[1]{}
Determine somente a contribuição da fonte de corrente para $V_A$.
\circuitoq{\begin{circuitikz}[american,scale=.84,transform shape]
\coordinate (g) at (3.4,0); \coordinate (a) at (3.4,3.6);
\draw (0,0) to[\fonteV,l^={$\SI{18}{\volt}$}] (0,3.6) to[R,l={$\SI{9}{\ohm}$}] (a);
\draw (a) to[R,l_={$\SI{6}{\ohm}$}] (g)--(0,0);
\draw (6.8,0) to[isource,l_={$\SI{3}{\ampere}$}] (6.8,3.6)--(a); \draw (6.8,0)--(g);
\node[above,Vcol] at(a){$V_A$}; \node[ground] at(g){};
\end{circuitikz}}
\resposta{$V_A^{(I)}=\SI{10.8}{\volt}$}
\resolucao{Com a fonte de tensão em curto, o nó vê $9\parallel6=\SI{3.6}{\ohm}$. Logo $V_A^{(I)}=3\times3{,}6=\SI{10.8}{\volt}$.}
\end{questao}

\begin{questao}[1]{}
Quantas análises parciais são necessárias para obter $V$ por superposição no circuito abaixo?
\circuitoq{\begin{circuitikz}[american,scale=.7,transform shape]
\coordinate (g) at (5,0); \coordinate (a) at (5,4);
\draw (0,0) to[\fonteV,l^={$E_1$}] (0,4) to[R,l={$R_1$}] (a);
\draw (2.5,0) to[isource,l^={$I_2$}] (2.5,4)--(a);
\draw (10,0) to[\fonteV,l_={$E_3$}] (10,4) to[R,l_={$R_3$}] (a);
\draw (a) to[R,l_={$R_L$}] (g); \draw (0,0)--(2.5,0)--(g)--(10,0);
\node[above,Vcol] at(a){$V$}; \node[ground] at(g){};
\end{circuitikz}}
\resposta{Três análises parciais}
\resolucao{Há três fontes independentes. Faz-se uma análise com $E_1$, outra com $I_2$ e outra com $E_3$, mantendo eventuais fontes dependentes ativas.}
\end{questao}

\begin{questao}[1]{}
As contribuições das duas fontes para a corrente indicada são $+\SI{1.8}{\ampere}$ e $-\SI{0.7}{\ampere}$. Determine a corrente real.
\circuitoq{\begin{circuitikz}[american,scale=.86,transform shape]
\draw (0,0) to[\fonteV,l^={$E_1$}] (0,3.4) to[R,l={$R$},i>^={$I$}] (4.8,3.4)
 to[\fonteV,l_={$E_2$}] (4.8,0)--(0,0);
\end{circuitikz}}
\resposta{$I=\SI{1.1}{\ampere}$ no sentido indicado}
\resolucao{Somam-se contribuições com sinal: $I=1{,}8-0{,}7=\SI{1.1}{\ampere}$.}
\end{questao}

\begin{questao}[1]{}
O teorema da superposição pode ser aplicado diretamente a este circuito para obter a tensão no diodo?
\circuitoq{\begin{circuitikz}[american,scale=.86,transform shape]
\draw (0,0) to[\fonteV,l^={$\SI{10}{\volt}$}] (0,3.5) to[R,l={$\SI{1}{\kilo\ohm}$}] (3.5,3.5)
 to[D,l={$D$}] (3.5,0)--(0,0);
\end{circuitikz}}
\resposta{Não}
\resolucao{O diodo é não linear: sua relação tensão-corrente não é proporcional. A configuração de condução pode mudar entre análises parciais, invalidando a soma.}
\end{questao}

% N2 --- Superposicao: 9 questoes, circuitos e valores distintos
\blocoaplicacao

\begin{questao}[2]{}
Determine $V_A$ por superposição.
\circuitoq{\begin{circuitikz}[american,scale=.76,transform shape]
\coordinate (g) at (5,0); \coordinate (a) at (5,4);
\draw (0,0) to[\fonteV,l^={$\SI{18}{\volt}$}] (0,4) to[R,l={$\SI{6}{\ohm}$}] (a);
\draw (10,0) to[\fonteV,l_={$\SI{6}{\volt}$}] (10,4) to[R,l_={$\SI{12}{\ohm}$}] (a);
\draw (a) to[R,l_={$\SI{4}{\ohm}$}] (g); \draw (0,0)--(g)--(10,0);
\node[above,Vcol] at(a){$V_A$}; \node[ground] at(g){};
\end{circuitikz}}
\resposta{$V_A=\SI{7}{\volt}$}
\resolucao{Com apenas $18$ V: $V_A'=\SI{6}{\volt}$. Com apenas $6$ V: $V_A''=\SI{1}{\volt}$. Logo $V_A=7$ V.}
\end{questao}

\begin{questao}[2]{}
Use superposição para obter a tensão sobre $R_L$.
\circuitoq{\begin{circuitikz}[american,scale=.84,transform shape]
\coordinate (g) at (3.8,0); \coordinate (a) at (3.8,3.8);
\draw (0,0) to[\fonteV,l^={$\SI{24}{\volt}$}] (0,3.8) to[R,l={$\SI{8}{\ohm}$}] (a);
\draw (a) to[R,l_={$R_L=\SI{12}{\ohm}$}] (g)--(0,0);
\draw (7.3,0) to[isource,l_={$\SI{1.5}{\ampere}$}] (7.3,3.8)--(a); \draw (7.3,0)--(g);
\node[above,Vcol] at(a){$V_A$}; \node[ground] at(g){};
\end{circuitikz}}
\resposta{$V_A=\SI{21.6}{\volt}$}
\resolucao{A resistência vista pelo nó com a fonte de tensão zerada é $8\parallel12=\SI{4.8}{\ohm}$. A fonte de tensão contribui com $14{,}4$ V e a de corrente com $7{,}2$ V; soma $21{,}6$ V.}
\end{questao}

\begin{questao}[2]{}
Determine, por superposição, a corrente no resistor que liga os nós 1 e 2, positiva de 1 para 2.
\circuitoq{\begin{circuitikz}[american,scale=.82,transform shape]
\coordinate (g) at (4,0); \coordinate (n1) at (2,3.8); \coordinate (n2) at (6.5,3.8);
\draw (0,0) to[isource,l^={$\SI{3}{\ampere}$}] (0,3.8)--(n1);
\draw (n1) to[R,l_={$\SI{6}{\ohm}$}] (2,0)--(g);
\draw (n1) to[R,l={$\SI{6}{\ohm}$},i>^={$i_{12}$}] (n2);
\draw (n2) to[R,l={$\SI{3}{\ohm}$}] (6.5,0)--(g);
\draw (8.5,0) to[isource,l_={$\SI{1}{\ampere}$}] (8.5,3.8)--(n2); \draw (8.5,0)--(g)--(0,0);
\node[above,Vcol] at(n1){$V_1$}; \node[above,Vcol] at(n2){$V_2$}; \node[ground] at(g){};
\end{circuitikz}}
\resposta{$V_1=\SI{12}{\volt}$, $V_2=\SI{6}{\volt}$, $i_{12}=\SI{1}{\ampere}$}
\resolucao{A soma das respostas parciais fornece $V_1=12$ V e $V_2=6$ V. Portanto $i_{12}=(12-6)/6=\SI{1}{\ampere}$.}
\end{questao}

\begin{questao}[2]{}
As fontes estão em oposição. Determine $I$ por superposição.
\circuitoq{\begin{circuitikz}[american,scale=.88,transform shape]
\draw (0,0) to[\fonteV,l^={$\SI{21}{\volt}$}] (0,3.4) to[R,l={$\SI{6}{\ohm}$},i>^={$I$}] (5,3.4)
 to[\fonteV,l_={$\SI{9}{\volt}$},invert] (5,0)--(0,0);
\end{circuitikz}}
\resposta{$I=\SI{2}{\ampere}$}
\resolucao{$I_1=21/6=3{,}5$ A e $I_2=-9/6=-1{,}5$ A. Assim $I=2$ A.}
\end{questao}

\begin{questao}[2]{}
Determine a corrente no resistor de $\SI{3}{\ohm}$ entre os nós usando a contribuição das três fontes separadamente.
\circuitoq{\begin{circuitikz}[american,scale=.68,transform shape]
\coordinate (g) at (5,0); \coordinate (n1) at (3.2,4.2); \coordinate (n2) at (7,4.2);
\draw (0,0) to[\fonteV,l^={$\SI{18}{\volt}$}] (0,4.2) to[R,l={$\SI{3}{\ohm}$}] (n1);
\draw (n1) to[R,l_={$\SI{6}{\ohm}$}] (3.2,0)--(g);
\draw (n1) to[R,l={$\SI{3}{\ohm}$},i>^={$i$}] (n2);
\draw (10.3,0) to[\fonteV,l_={$\SI{6}{\volt}$}] (10.3,4.2) to[R,l_={$\SI{6}{\ohm}$}] (n2);
\draw (n2) to[R,l={$\SI{3}{\ohm}$}] (7,0)--(g);
\draw (8.7,0) to[isource,l_={$\SI{1}{\ampere}$}] (8.7,4.2)--(n2); \draw (0,0)--(g)--(8.7,0)--(10.3,0);
\end{circuitikz}}
\resposta{$V_1=\SI{68/7}{\volt}$, $V_2=\SI{44/7}{\volt}$, $i=\SI{8/7}{\ampere}$}
\resolucao{Somando as três respostas parciais obtêm-se $V_1=68/7$ V e $V_2=44/7$ V. Logo $i=(V_1-V_2)/3=8/7$ A.}
\end{questao}

\begin{questao}[2]{}
As duas fontes alimentam uma ponte resistiva. Determine por superposição a tensão $V_{ab}=V_a-V_b$.
\circuitoq{\begin{circuitikz}[american,scale=.72,transform shape]
\coordinate (l) at (0,2.2); \coordinate (r) at (10,2.2); \coordinate (a) at (5,5); \coordinate (b) at (5,-.6);
\draw (l) to[\fonteV,l^={$\SI{12}{\volt}$}] (0,-.6)--(b);
\draw (r) to[\fonteV,l_={$\SI{6}{\volt}$},invert] (10,-.6)--(b);
\draw (l) to[R,l={$\SI{4}{\ohm}$}] (a) to[R,l={$\SI{8}{\ohm}$}] (r);
\draw (l) to[R,l_={$\SI{6}{\ohm}$}] (b) to[R,l_={$\SI{12}{\ohm}$}] (r);
\node[above,Vcol] at(a){$a$}; \node[below,Vcol] at(b){$b$};
\end{circuitikz}}
\resposta{$V_{ab}=\SI{0}{\volt}$}
\resolucao{Os dois divisores têm a mesma razão $4/8=6/12$. Para qualquer combinação linear das duas fontes, os pontos médios permanecem equipotenciais; cada fonte separadamente produz $V_{ab}=0$, portanto a soma também é zero.}
\end{questao}

\begin{questao}[2]{}
Uma fonte real e uma fonte ideal de corrente alimentam a carga. Determine $V$ por superposição.
\circuitoq{\begin{circuitikz}[american,scale=.82,transform shape]
\coordinate (g) at (3.8,0); \coordinate (a) at (3.8,3.8);
\draw (0,0) to[\fonteV,l^={$\SI{15}{\volt}$}] (0,3.8) to[R,l={$r=\SI{3}{\ohm}$}] (a);
\draw (a) to[R,l_={$R_L=\SI{6}{\ohm}$}] (g)--(0,0);
\draw (7.2,0) to[isource,l_={$\SI{2}{\ampere}$}] (7.2,3.8)--(a); \draw (7.2,0)--(g);
\node[above,Vcol] at(a){$V$}; \node[ground] at(g){};
\end{circuitikz}}
\resposta{$V=\SI{14}{\volt}$}
\resolucao{A fonte real contribui com $10$ V. Com a f.e.m. zerada, $3\parallel6=2\,\Omega$ e a fonte de corrente contribui com $4$ V. Total: $14$ V.}
\end{questao}

\begin{questao}[2]{}
Determine a contribuição de cada fonte e a tensão na carga.
\circuitoq{\begin{circuitikz}[american,scale=.76,transform shape]
\coordinate (g) at (5,0); \coordinate (a) at (5,4);
\draw (0,0) to[\fonteV,l^={$\SI{20}{\volt}$}] (0,4) to[R,l={$\SI{4}{\ohm}$}] (a);
\draw (10,0) to[\fonteV,l_={$\SI{12}{\volt}$}] (10,4) to[R,l_={$\SI{12}{\ohm}$}] (a);
\draw (a) to[R,l_={$R_L=\SI{6}{\ohm}$}] (g); \draw (0,0)--(g)--(10,0);
\node[above,Vcol] at(a){$V_L$}; \node[ground] at(g){};
\end{circuitikz}}
\resposta{$V_L^{(20)}=\SI{10}{\volt}$, $V_L^{(12)}=\SI{2}{\volt}$, $V_L=\SI{12}{\volt}$}
\resolucao{Com a fonte de 12 V em curto, $4$ alimenta $12\parallel6=4\,\Omega$, dando $10$ V. Com a de 20 V em curto, $12$ alimenta $4\parallel6=2{,}4\,\Omega$, dando $2$ V. Soma: $12$ V.}
\end{questao}

\begin{questao}[2]{}
Use superposição para determinar $V_2$.
\circuitoq{\begin{circuitikz}[american,scale=.76,transform shape]
\coordinate (g) at (4.8,0); \coordinate (n1) at (3,4); \coordinate (n2) at (7,4);
\draw (0,0) to[\fonteV,l^={$\SI{24}{\volt}$}] (0,4) to[R,l={$\SI{6}{\ohm}$}] (n1);
\draw (n1) to[R,l_={$\SI{6}{\ohm}$}] (3,0)--(g);
\draw (n1) to[R,l={$\SI{3}{\ohm}$}] (n2);
\draw (n2) to[R,l={$\SI{4}{\ohm}$}] (7,0)--(g);
\draw (9.2,0) to[isource,l_={$\SI{1}{\ampere}$}] (9.2,4)--(n2); \draw (9.2,0)--(g)--(0,0);
\node[above,Vcol] at(n1){$V_1$}; \node[above,Vcol] at(n2){$V_2$}; \node[ground] at(g){};
\end{circuitikz}}
\resposta{$V_2=\SI{7.2}{\volt}$}
\resolucao{A soma das respostas parciais da fonte de 24 V e da fonte de 1 A fornece $V_1=9{,}6$ V e $V_2=7{,}2$ V.}
\end{questao}

% N3 --- Superposicao: 10 questoes com topologias distintas
% Revisao visual: ramos mais espaçados, rotulos afastados e sem diagonais
% desnecessarias ou sobreposicao entre componentes.
\blocoaprofundamento

\begin{questao}[3]{}
A fonte dependente deve permanecer ativa em todas as parcelas. Determine $V_A$ por superposição das duas fontes independentes.
\circuitoq{\begin{circuitikz}[american,scale=.74,transform shape,line width=.8pt]
\coordinate (g) at (4.5,0); \coordinate (a) at (4.5,4.4);
\draw (0,0) to[\fonteV,l^={$\SI{12}{\volt}$}] (0,4.4)
      to[R,l={$\SI{4}{\ohm}$}] (a);
\draw (a) to[R,l_={$\SI{8}{\ohm}$}] (g)--(0,0);
\draw (8.0,0) to[isource,l_={$\SI{1}{\ampere}$}] (8.0,4.4)--(a);
\draw (11.6,0) to[cisource,l_={$0{,}0625V_A$}] (11.6,4.4)--(8.0,4.4);
\draw (11.6,0)--(8.0,0)--(g);
\fill (a) circle (1.1pt);
\node[Vcol,fill=white,inner sep=1.6pt] at ($(a)+(0,.68)$) {$V_A$};
\node[ground] at(g){};
\end{circuitikz}}
\resposta{$V_A'=\SI{9.6}{\volt}$, $V_A''=\SI{3.2}{\volt}$, $V_A=\SI{12.8}{\volt}$}
\resolucao{A condutância líquida é $1/4+1/8-0{,}0625=0{,}3125$ S. A fonte de 12 V injeta equivalente de 3 A, dando $9{,}6$ V; a fonte de 1 A dá $3{,}2$ V. A dependente permanece nas duas análises.}
\end{questao}

\begin{questao}[3]{}
Há três fontes independentes atuando em dois nós. Determine $V_1$ e $V_2$ por superposição.
\circuitoq{\begin{circuitikz}[american,scale=.62,transform shape,line width=.8pt]
\coordinate (g) at (6.2,0); \coordinate (n1) at (4.1,4.8); \coordinate (n2) at (8.5,4.8);
\draw (0,0) to[\fonteV,l^={$\SI{18}{\volt}$}] (0,4.8)
      to[R,l={$\SI{6}{\ohm}$}] (n1);
\draw (2.0,0) to[isource,l^={$\SI{2}{\ampere}$}] (2.0,4.8)--(n1);
\draw (n1) to[R,l_={$\SI{9}{\ohm}$}] (4.1,0)--(g);
\draw (n1) to[R,l={$\SI{3}{\ohm}$}] (n2);
\draw (n2) to[R,l={$\SI{6}{\ohm}$}] (8.5,0)--(g);
\draw (12.8,0) to[\fonteV,l_={$\SI{9}{\volt}$}] (12.8,4.8)
      to[R,l_={$\SI{3}{\ohm}$}] (n2);
\draw (0,0)--(2.0,0)--(g)--(12.8,0);
\fill (n1) circle (1.1pt); \fill (n2) circle (1.1pt);
\node[Vcol,fill=white,inner sep=1.6pt] at ($(n1)+(0,.72)$) {$V_1$};
\node[Vcol,fill=white,inner sep=1.6pt] at ($(n2)+(0,.72)$) {$V_2$};
\node[ground] at(g){};
\end{circuitikz}}
\resposta{$V_1\approx\SI{12.98}{\volt}$; $V_2\approx\SI{8.79}{\volt}$}
\resolucao{Resolvem-se três vezes as mesmas duas equações nodais, uma por fonte. A soma das parcelas fornece $V_1=558/43\approx12{,}98$ V e $V_2=378/43\approx8{,}79$ V.}
\end{questao}

\begin{questao}[3]{}
Duas fontes reais alimentam uma carga. Determine a tensão $U$ e identifique a parcela de cada f.e.m.
\circuitoq{\begin{circuitikz}[american,scale=.74,transform shape,line width=.8pt]
\coordinate (g) at (5.3,0); \coordinate (a) at (5.3,4.4);
\draw (0,0) to[\fonteV,l^={$\SI{18}{\volt}$}] (0,4.4)
      to[R,l={$r_1=\SI{3}{\ohm}$}] (a);
\draw (11.0,0) to[\fonteV,l_={$\SI{9}{\volt}$}] (11.0,4.4)
      to[R,l_={$r_2=\SI{6}{\ohm}$}] (a);
\draw (a) to[R,l_={$R_L=\SI{9}{\ohm}$}] (g);
\draw (0,0)--(g)--(11.0,0);
\fill (a) circle (1.1pt);
\node[Vcol,fill=white,inner sep=1.6pt] at ($(a)+(0,.68)$) {$U$};
\node[ground] at(g){};
\end{circuitikz}}
\resposta{$U_1\approx\SI{9.82}{\volt}$; $U_2\approx\SI{2.45}{\volt}$; $U\approx\SI{12.27}{\volt}$}
\resolucao{A condutância total é $1/3+1/6+1/9=11/18$ S. As parcelas são $(18/3)/(11/18)=108/11$ V e $(9/6)/(11/18)=27/11$ V. Soma $135/11\approx12{,}27$ V.}
\end{questao}

\begin{questao}[3]{}
A fonte de tensão controlada impõe $V_2=1{,}5V_1$. Determine $V_1$ por superposição das fontes independentes.
\circuitoq{\begin{circuitikz}[american,scale=.72,transform shape,line width=.8pt]
\coordinate (g) at (3.0,0); \coordinate (n1) at (3.0,4.5); \coordinate (n2) at (8.2,4.5);
\draw (0,0) to[isource,l^={$\SI{2}{\ampere}$}] (0,4.5)--(n1);
\draw (n1) to[R,l_={$\SI{6}{\ohm}$}] (g);
\draw (n1) to[R,l={$\SI{12}{\ohm}$}] (n2);
\draw (n2) to[cvsource,l_={$1{,}5V_1$}] (8.2,0)--(g)--(0,0);
\draw (11.2,0) to[isource,l_={$\SI{0.5}{\ampere}$}] (11.2,4.5)--(n2);
\draw (11.2,0)--(8.2,0);
\fill (n1) circle (1.1pt); \fill (n2) circle (1.1pt);
\node[Vcol,fill=white,inner sep=1.6pt] at ($(n1)+(0,.72)$) {$V_1$};
\node[Vcol,fill=white,inner sep=1.6pt] at ($(n2)+(0,.72)$) {$V_2$};
\node[ground] at(g){};
\end{circuitikz}}
\resposta{A fonte dependente permanece ativa em ambas as análises; $V_1=\SI{16}{\volt}$}
\resolucao{A superposição é feita somente sobre as duas fontes independentes; a fonte controlada não é zerada. Montando as equações em cada parcela e somando, obtém-se $V_1=16$ V e $V_2=24$ V.}
\end{questao}

\begin{questao}[3]{}
Mostre quantitativamente por que não se deve somar potências parciais no resistor.
\circuitoq{\begin{circuitikz}[american,scale=.90,transform shape,line width=.8pt]
\draw (0,0) to[\fonteV,l^={$\SI{16}{\volt}$}] (0,3.8)
      to[R,l={$R=\SI{4}{\ohm}$},i>^={$I$}] (5.6,3.8)
      to[\fonteV,l_={$\SI{8}{\volt}$}] (5.6,0)--(0,0);
\end{circuitikz}}
\resposta{$I_1=\SI{4}{\ampere}$, $I_2=\SI{2}{\ampere}$; $P=\SI{144}{\watt}$, enquanto $P_1+P_2=\SI{80}{\watt}$}
\resolucao{As correntes somam: $I=4+2=6$ A. Portanto $P=RI^2=4(36)=144$ W. Somar $4(4^2)+4(2^2)=80$ W ignora o termo cruzado $2RI_1I_2=64$ W.}
\end{questao}

\begin{questao}[3]{}
Use o circuito para explicar por que a superposição falha quando o estado do diodo muda entre parcelas.
\circuitoq{\begin{circuitikz}[american,scale=.74,transform shape,line width=.8pt]
\coordinate (g) at (5.0,0); \coordinate (a) at (5.0,4.4);
\draw (0,0) to[\fonteV,l^={$\SI{6}{\volt}$}] (0,4.4)
      to[R,l={$\SI{1}{\kilo\ohm}$}] (a);
\draw (10.4,0) to[\fonteV,l_={$\SI{4}{\volt}$},invert] (10.4,4.4)
      to[R,l_={$\SI{2}{\kilo\ohm}$}] (a);
\draw (a) to[D,l_={$D$}] (g);
\draw (0,0)--(g)--(10.4,0);
\fill (a) circle (1.1pt); \node[ground] at(g){};
\end{circuitikz}}
\resposta{Não se pode superpor as respostas obtidas com estados diferentes do diodo}
\resolucao{O diodo altera a própria topologia equivalente ao conduzir ou bloquear. Assim, cada parcela pode corresponder a uma rede diferente; a linearidade exigida pelo teorema deixa de existir.}
\end{questao}

\begin{questao}[3]{}
No circuito linear, determine o novo $V_A$ se a fonte de $\SI{4}{\volt}$ for duplicada, sem refazer toda a análise.
\circuitoq{\begin{circuitikz}[american,scale=.76,transform shape,line width=.8pt]
\coordinate (g) at (5.2,0); \coordinate (a) at (5.2,4.3);
\draw (0,0) to[\fonteV,l^={$\SI{10}{\volt}$}] (0,4.3)
      to[R,l={$\SI{5}{\ohm}$}] (a);
\draw (10.8,0) to[\fonteV,l_={$\SI{4}{\volt}$}] (10.8,4.3)
      to[R,l_={$\SI{10}{\ohm}$}] (a);
\draw (a) to[R,l_={$\SI{10}{\ohm}$}] (g);
\draw (0,0)--(g)--(10.8,0);
\fill (a) circle (1.1pt);
\node[Vcol,fill=white,inner sep=1.6pt] at ($(a)+(0,.68)$) {$V_A$};
\node[ground] at(g){};
\end{circuitikz}}
\resposta{Originalmente $V_A=\SI{6}{\volt}$; após duplicar a segunda fonte, $V_A=\SI{7}{\volt}$}
\resolucao{As parcelas originais são $5$ V e $1$ V. Duplicar apenas a segunda fonte duplica somente sua parcela: $1\to2$ V. Logo o novo total é $5+2=7$ V.}
\end{questao}

\begin{questao}[3]{}
A corrente total no galvanômetro é nula. A fonte $E_1$, sozinha, produz $+\SI{35}{\milli\ampere}$. Determine a contribuição de $E_2$ e interprete.
\circuitoq{\begin{circuitikz}[american,scale=.62,transform shape,line width=.8pt]
\coordinate (l) at (0,2.5); \coordinate (r) at (12.0,2.5);
\coordinate (u) at (6.0,6.2); \coordinate (d) at (6.0,-1.2);
\draw (l) to[\fonteV,l^={$E_1$}] (0,-1.2)--(d);
\draw (r) to[\fonteV,l_={$E_2$}] (12.0,-1.2)--(d);
\draw (l) to[R,l={$R_1$}] (u) to[R,l={$R_3$}] (r);
\draw (l) to[R,l_={$R_2$}] (d) to[R,l_={$R_4$}] (r);
\draw (u) to[R,l={$R_g$},i>^={$I_g$}] (d);
\node[Vcol,fill=white,inner sep=1.5pt] at ($(u)+(0,.55)$) {$a$};
\node[Vcol,fill=white,inner sep=1.5pt] at ($(d)+(0,-.55)$) {$b$};
\node[ground] at(d){};
\end{circuitikz}}
\resposta{$I_g^{(E_2)}=\SI{-35}{\milli\ampere}$}
\resolucao{Como $I_g=I_g^{(E_1)}+I_g^{(E_2)}=0$, a segunda parcela deve ser $-35$ mA. Isso é cancelamento entre fontes, não necessariamente equilíbrio resistivo da ponte.}
\end{questao}

\begin{questao}[3]{}
Um sinal e uma interferência entram por ramos diferentes. Determine as parcelas de saída e a razão sinal-interferência.
\circuitoq{\begin{circuitikz}[american,scale=.72,transform shape,line width=.8pt]
\coordinate (g) at (5.3,0); \coordinate (a) at (5.3,4.5);
\draw (0,0) to[\fonteV,l^={$V_s=\SI{0.2}{\volt}$}] (0,4.5)
      to[R,l={$\SI{1}{\kilo\ohm}$}] (a);
\draw (11.0,0) to[\fonteV,l_={$V_n=\SI{1}{\volt}$}] (11.0,4.5)
      to[R,l_={$\SI{20}{\kilo\ohm}$}] (a);
\draw (a) to[R,l_={$R_L=\SI{2}{\kilo\ohm}$}] (g);
\draw (0,0)--(g)--(11.0,0);
\fill (a) circle (1.1pt);
\node[Vcol,fill=white,inner sep=1.6pt] at ($(a)+(0,.70)$) {$V_o$};
\node[ground] at(g){};
\end{circuitikz}}
\resposta{$V_{o,s}\approx\SI{129}{\milli\volt}$; $V_{o,n}\approx\SI{32.3}{\milli\volt}$; razão $\approx4$ ($\approx\SI{12}{\decibel}$)}
\resolucao{A condutância total é $1/1k+1/20k+1/2k=1{,}55$ mS. As injeções são $0{,}2$ mA e $0{,}05$ mA, gerando aproximadamente $129$ mV e $32{,}3$ mV.}
\end{questao}

\begin{questao}[3]{}
Compare o custo de superposição e análise nodal para a rede de dois nós e quatro fontes abaixo.
\circuitoq{\begin{circuitikz}[american,scale=.56,transform shape,line width=.8pt]
\coordinate (g) at (7.0,0); \coordinate (n1) at (4.4,5.0); \coordinate (n2) at (9.4,5.0);
\draw (0,0) to[\fonteV,l^={$E_1$}] (0,5.0) to[R,l={$R_1$}] (n1);
\draw (2.1,0) to[isource,l^={$I_2$}] (2.1,5.0)--(n1);
\draw (n1) to[R,l_={$R_2$}] (4.4,0)--(g);
\draw (n1) to[R,l={$R_{12}$}] (n2);
\draw (n2) to[R,l={$R_3$}] (9.4,0)--(g);
\draw (11.8,0) to[isource,l_={$I_4$}] (11.8,5.0)--(n2);
\draw (14.0,0) to[\fonteV,l_={$E_3$}] (14.0,5.0) to[R,l_={$R_4$}] (n2);
\draw (0,0)--(2.1,0)--(g)--(11.8,0)--(14.0,0);
\fill (n1) circle (1.1pt); \fill (n2) circle (1.1pt);
\node[Vcol,fill=white,inner sep=1.5pt] at ($(n1)+(0,.72)$) {$V_1$};
\node[Vcol,fill=white,inner sep=1.5pt] at ($(n2)+(0,.72)$) {$V_2$};
\node[ground] at(g){};
\end{circuitikz}}
\resposta{Superposição: quatro sistemas $2\times2$; nodal direto: um sistema $2\times2$}
\resolucao{Para um único resultado numérico, nodal é mais econômico. Superposição vale o custo adicional quando é necessário conhecer a contribuição individual de cada fonte.}
\end{questao}
% N4 --- Superposicao: 8 desafios, com circuitos quando pedagogicamente util
% Revisao visual preventiva: maior separacao entre ramos, fontes e rotulos.
\blocodesafio

\begin{questao}[4]{}
A rede linear abaixo é vista por um nó de saída. Demonstre, usando a forma matricial nodal, que a resposta total é a soma das respostas às fontes $E_1$, $I_2$ e $E_3$ separadamente.
\circuitoq{\begin{circuitikz}[american,scale=.68,transform shape,line width=.8pt]
\coordinate (g) at (5.8,0); \coordinate (a) at (5.8,4.5);
\draw (0,0) to[\fonteV,l^={$E_1$}] (0,4.5) to[R,l={$R_1$}] (a);
\draw (8.8,0) to[isource,l_={$I_2$}] (8.8,4.5)--(a);
\draw (12.0,0) to[\fonteV,l_={$E_3$}] (12.0,4.5) to[R,l_={$R_3$}] (a);
\draw (a) to[R,l_={$R_L$}] (g);
\draw (0,0)--(g)--(8.8,0)--(12.0,0);
\fill (a) circle (1.1pt); \node[Vcol,fill=white,inner sep=1.5pt] at ($(a)+(0,.68)$) {$V_o$};
\node[ground] at(g){};
\end{circuitikz}}
\resposta{$\mathbf v=\mathbf G^{-1}\mathbf b_1+\mathbf G^{-1}\mathbf b_2+\mathbf G^{-1}\mathbf b_3$}
\resolucao{Em uma rede linear, $\mathbf G$ não muda com o valor das fontes e $\mathbf b=\mathbf b_1+\mathbf b_2+\mathbf b_3$. Logo $\mathbf v=\mathbf G^{-1}\mathbf b=\sum_k\mathbf G^{-1}\mathbf b_k$. Cada termo é exatamente uma análise parcial.}
\end{questao}

\begin{questao}[4]{}
No resistor, as parcelas de corrente são $I_1$ e $I_2$. Expanda a potência e determine a condição para que a soma das potências parciais coincida com a potência real.
\circuitoq{\begin{circuitikz}[american,scale=.88,transform shape,line width=.8pt]
\coordinate (g) at (3.8,0); \coordinate (a) at (3.8,4.0);
\draw (0,0) to[isource,l^={$I_1$}] (0,4.0)--(a);
\draw (7.6,0) to[isource,l_={$I_2$}] (7.6,4.0)--(a);
\draw (a) to[R,l={$R$},i>^={$I$}] (g);
\draw (0,0)--(g)--(7.6,0);
\end{circuitikz}}
\resposta{$P=RI_1^2+RI_2^2+2RI_1I_2$; a soma parcial só coincide quando $I_1I_2=0$}
\resolucao{O termo cruzado $2RI_1I_2$ é a parcela ausente na soma ingênua. Ele pode ser positivo ou negativo conforme as contribuições tenham o mesmo sentido ou sentidos opostos.}
\end{questao}

\begin{questao}[4]{}
Projete um ensaio experimental de superposição no circuito, especificando as três configurações de medição e o modo correto de desativar cada fonte.
\circuitoq{\begin{circuitikz}[american,scale=.80,transform shape,line width=.8pt]
\coordinate (g) at (4.5,0); \coordinate (a) at (4.5,4.2);
\draw (0,0) to[\fonteV,l^={$E=\SI{15}{\volt}$}] (0,4.2) to[R,l={$\SI{5}{\ohm}$}] (a);
\draw (a) to[R,l_={$\SI{10}{\ohm}$}] (g)--(0,0);
\draw (8.8,0) to[isource,l_={$\SI{1}{\ampere}$}] (8.8,4.2)--(a); \draw (8.8,0)--(g);
\fill (a) circle (1.1pt); \node[Vcol,fill=white,inner sep=1.6pt] at ($(a)+(0,.68)$) {$V_A$};
\node[ground] at(g){};
\end{circuitikz}}
\resposta{Medir com ambas ativas, depois somente $E$, depois somente $I$; verificar $V_A\approx V_A^{(E)}+V_A^{(I)}$}
\resolucao{Ao medir apenas $E$, abre-se o ramo da fonte de corrente. Ao medir apenas $I$, substitui-se a fonte de tensão ideal por curto. Não se deve simplesmente retirar a fonte de tensão do circuito.}
\end{questao}

\begin{questao}[4]{}
Escolha $E_2$ para anular a corrente em $R$ sem alterar $E_1$.
\circuitoq{\begin{circuitikz}[american,scale=.90,transform shape,line width=.8pt]
\draw (0,0) to[\fonteV,l^={$E_1=\SI{30}{\volt}$}] (0,3.8)
      to[R,l={$R=\SI{10}{\ohm}$},i>^={$I$}] (5.8,3.8)
      to[\fonteV,l_={$E_2$},invert] (5.8,0)--(0,0);
\end{circuitikz}}
\resposta{$E_2=\SI{30}{\volt}$ com polaridade oposta à de $E_1$}
\resolucao{A parcela de $E_1$ é $+3$ A. Para corrente total nula, a segunda deve produzir $-3$ A; portanto $E_2=10(3)=30$ V em oposição.}
\end{questao}

\begin{questao}[4]{}
Com as fontes independentes zeradas, determine a resistência vista na porta aplicando uma fonte de teste, mantendo a fonte dependente ativa.
\circuitoq{\begin{circuitikz}[american,scale=.76,transform shape,line width=.8pt]
\coordinate (g) at (4.2,0); \coordinate (a) at (4.2,4.3);
\draw (a) to[R,l_={$\SI{1}{\kilo\ohm}$}] (g);
\draw (0,0) to[cisource,l^={$0{,}5\,\mathrm{mS}\,V_A$}] (0,4.3)--(a); \draw (0,0)--(g);
\draw (8.8,0) to[isource,l_={$I_t$}] (8.8,4.3)--(a); \draw (8.8,0)--(g);
\fill (a) circle (1.1pt); \node[Vcol,fill=white,inner sep=1.6pt] at ($(a)+(0,.68)$) {$V_A$};
\node[ground] at(g){};
\end{circuitikz}}
\resposta{$R_{eq}=\SI{2}{\kilo\ohm}$}
\resolucao{$I_t=V_A/1000-0{,}0005V_A=0{,}0005V_A$. Assim $V_A/I_t=2000\,\Omega$. A dependente não pode ser desativada.}
\end{questao}

\begin{questao}[4]{}
Uma rede tem quatro fontes, mas deseja-se apenas saber qual delas domina $V_o$. Proponha uma estratégia de superposição que minimize o trabalho e produza uma medida quantitativa de dominância.
\circuitoq{\begin{circuitikz}[american,scale=.58,transform shape,line width=.8pt]
\coordinate (g) at (7.0,0); \coordinate (a) at (7.0,5.0);
\draw (0,0) to[\fonteV,l^={$E_1$}] (0,5.0) to[R,l={$R_1$}] (a);
\draw (2.5,0) to[isource,l^={$I_2$}] (2.5,5.0)--(a);
\draw (11.5,0) to[isource,l_={$I_4$}] (11.5,5.0)--(a);
\draw (14.0,0) to[\fonteV,l_={$E_3$}] (14.0,5.0) to[R,l_={$R_3$}] (a);
\draw (a) to[R,l_={$R_L$}] (g);
\draw (0,0)--(2.5,0)--(g)--(11.5,0)--(14.0,0);
\fill (a) circle (1.1pt); \node[Vcol,fill=white,inner sep=1.5pt] at ($(a)+(0,.72)$) {$V_o$};
\node[ground] at(g){};
\end{circuitikz}}
\resposta{Calcular as quatro parcelas $V_{o,k}$ e comparar $|V_{o,k}|/\sum_j|V_{o,j}|$}
\resolucao{A fatoração da matriz nodal pode ser reutilizada nas quatro excitações. O índice normalizado pelos módulos evita que cancelamentos algébricos escondam uma fonte individualmente forte.}
\end{questao}

\begin{questao}[4]{}
Formule um algoritmo para calcular simultaneamente $n$ saídas de uma rede linear com $m$ fontes, evitando refatorar a matriz a cada parcela.
\resposta{Fatorar $\mathbf G$ uma vez; resolver $m$ vetores de excitação; armazenar as respostas em uma matriz de transferência $\mathbf H$}
\resolucao{Monte $\mathbf G=\mathbf L\mathbf U$ uma única vez. Para cada fonte $k$, resolva $\mathbf G\mathbf v_k=\mathbf b_k$ por substituição. As $n$ grandezas selecionadas formam a coluna $k$ de $\mathbf H$, permitindo depois $\mathbf y=\mathbf H\mathbf s$.}
\end{questao}

\begin{questao}[4]{}
No circuito, indique até que faixa de operação a superposição é válida e qual evento físico faz o modelo deixar de ser linear.
\circuitoq{\begin{circuitikz}[american,scale=.74,transform shape,line width=.8pt]
\coordinate (g) at (4.4,0); \coordinate (a) at (4.4,4.2);
\draw (0,0) to[\fonteV,l^={$E_1$}] (0,4.2) to[R,l={$R_1$}] (a);
\draw (8.8,0) to[\fonteV,l_={$E_2$}] (8.8,4.2) to[R,l_={$R_2$}] (a);
\draw (a) to[R,l_={$R_L$}] (g); \draw (0,0)--(g)--(8.8,0);
\fill (a) circle (1.1pt); \node[Vcol,fill=white,inner sep=1.5pt] at ($(a)+(0,.68)$) {$V_o$};
\node[draw=cinza,rounded corners,align=center,font=\footnotesize,inner sep=3pt]
      at (11.4,2.1) {fonte real com\\limitação de corrente};
\node[ground] at(g){};
\end{circuitikz}}
\resposta{Válida enquanto resistores e fontes operarem em região linear; falha quando uma fonte entra em limitação/saturação ou um elemento muda de estado}
\resolucao{A superposição não é uma regra de desenho, mas consequência da linearidade do modelo. Se a fonte satura, a relação entre excitação e resposta deixa de ser proporcional e as parcelas já não podem ser somadas.}
\end{questao}
